\par\Needspace{6\baselineskip}\medskip
\begin{center}\hypertarget{native-179-article133}{}\pdfbookmark[2]{133.}{article-133}133.\end{center}
Da sich die Kosinus der Winkel $f,f',f''$ ebenso wie die Größen $\eta,\eta''$ um eine Größe zweiter Ordnung
von eins unterscheiden, begeht man einen Fehler vierter Ordnung, wenn man für $\dfrac{n+n''}{n'}$ den
Näherungswert $1+\dfrac{\theta\theta''}{2rr'r''}$ einsetzt. Verwendet man daher anstelle der Gleichung in Artikel 114 die
Gleichung
\[
a\delta'=b+\frac{c\theta+d\theta''}{\theta'}\left(1+\frac{\theta\theta''}{2rr'r''}\right)\mathpunct{\raisebox{.45ex}{,}}
\]
\hypertarget{native-180}{}so ergibt sich im Wert des Abstands $\delta'$ ein Fehler zweiter Ordnung, wenn die äußeren
Beobachtungen von der mittleren gleich weit entfernt sind, und in den übrigen Fällen ein Fehler
erster Ordnung. Diese neue Form der Gleichung eignet sich jedoch nicht zur Bestimmung von $\delta'$,
weil sie die noch unbekannten Größen $r,r',r''$ enthält.
\par
Im Allgemeinen unterscheiden sich die Größen $\dfrac{r}{r'}\mathpunct{\raisebox{.45ex}{,}}\;\dfrac{r''}{r'}$ ebenso wie das Produkt $\dfrac{rr''}{r'r'}$ um eine Größe
erster Ordnung von eins. In dem mehrfach erwähnten Sonderfall sieht man leicht, dass sich dieses
Produkt von eins nur um eine Größe zweiter Ordnung unterscheidet. Ist die Bahn eine Ellipse mit
geringer Exzentrizität, sodass man die Exzentrizität als Größe erster Ordnung betrachten darf,
so lässt sich der Unterschied zwischen $\dfrac{rr''}{r'r'}$ und eins sogar einer nochmals um eine Stufe
höheren Ordnung zuordnen. Der Fehler bleibt also offensichtlich von derselben Ordnung wie zuvor,
wenn man in unserer Gleichung für $\dfrac{\theta\theta''}{2rr'r''}$ den Ausdruck $\dfrac{\theta\theta''}{2r'^3}\mathpunct{\raisebox{.45ex}{,}}$ einsetzt. Die Gleichung nimmt damit
folgende Form an
\[
a\delta'=b+\frac{c\theta+d\theta''}{\theta'}\left(1+\frac{\theta\theta''}{2r'^3}\right)\mathpunct{\cdot}
\]
Diese Gleichung enthält zwar noch die unbekannte Größe $r'$, doch lässt sich diese offenbar
eliminieren, da sie nur von $\delta'$ und von bekannten Größen abhängt. Brächte man die Gleichung
anschließend in die übliche Form, so wäre sie vom \textit{achten} Grad.
